% Candidate: complete actual-field proof; no promotion before two full audits.
\subsection{A selected high point in the full residual-wild Iwahori image}

Let $K$ be complete discretely valued of mixed characteristic $(0,5)$,
with algebraically closed residue field, $e_K=20$, and $\mu _5\subset K$.
Suppose two whole potentially supersingular height-two groups have an
actual $125$-torsion isomorphism and different finite special characters.
Use the common finite twist and determinant normalization of B77 and
B91. In this subsection their residual wild action is nontrivial.
B91--B95 leave only the full Iwahori image of order $5^9$, the common
tame good field $L/K$ of degree24, and the both-large parameters
$B_i=5b_i>20$. Write $j_i=24-b_i$, so the actual high quotients
$Q_i=M_i/H$ have parameters $5j_i>20$. Their intrinsic old images
$M_i[5]/H$ are their common canonical high line.

All field valuations below have value group $\mathbf Z$. A conductor
means the last upper ramification break, without adding one. Fields
are generated by actual indicated points, and a statement about an
individual point of $Q_i[125]$ never asserts a comparison of all
$Q_i[125]$.

\begin{lemma}[The selected high point and its actual fields]
\label{lem:B-RD96-selected}
\status{Candidate; complete proof awaiting independent audits.}
Either the original parameter pair $(b_1,b_2)$ or its sorted opposite
pair $(j_1,j_2)$ can be used as an orientation. In each orientation
there is a common selected primitive high $125$-point $P^+$, a common
high $25$-point $P=5P^+$, and actual fields
\[
 E_h=L(P),\qquad S=R_oE_h,\qquad T=R_oL(P^+),
 \qquad [E_h:L]=25,\quad[S:L]=125,\quad[T:L]=3125,
                                                        \tag{B-RD96.1}
\]
where $R_o/L$ is the residual field of the chosen orientation.
The extension $T/S$ is Galois with group $\mathbf F_5^2$ and a
distinguished high-translation line $H_T$.

If the chosen parameters are $b_1,b_2$, its two coordinate functions
on $T$ have valuations and displacements
\[
\begin{gathered}
 A_i=25(24-b_i),\qquad
 D_{0,i}=3125b_i,\qquad D_{1,i}=15625(24-b_i),\\
 v_T(gP_i^+-P_i^+)=D_{1,i}\quad(0\ne g\in H_T),\qquad
 v_T(gP_i^+-P_i^+)=D_{0,i}\quad(g\notin H_T),\qquad
 v_T(5)=1500000.
                                                        \tag{B-RD96.2}
\end{gathered}
\]
The notation for displacement denotes the difference of formal
coordinates; its valuation equals that of the translating torsion
point.
\end{lemma}
\begin{proof}
Choose a common integral tame eigenbasis whose original high line is
$\mathbf F_5e_1$. The actual common division field $F/L$ has group
\[
 W=\left\{\begin{pmatrix}1+5a&v\\5w&1+5d\end{pmatrix}:
 a,d,w\in\mathbf Z/25,\ v\in\mathbf Z/125\right\}.
                                                        \tag{B-RD96.3}
\]
Indeed the actual image is contained in this group and B95 gives
its exact order $5^9$, equal to the order of the displayed group.
Write $a=a_0+5a_1$, $d=d_0+5d_1$, $w=w_0+5w_1$ and
$v=v_0+5v_1+25v_2$, with digits in $\mathbf F_5$.
There is no determinant relation between these digits.

For the original orientation let $P^+=e_1/125$ and $P=e_1/25$.
Their stabilizers are respectively $a=w=0$ and $a_0=w_0=0$.
The residual field $R_o$ is the cyclic field of $v_0$.
Consequently $S$ is fixed by $v_0=w_0=a_0=0$, and $T$ by
$v_0=w_0=a_0=a_1=w_1=0$. These stabilizers give all degrees
in (B-RD96.1). On the subgroup for $S$, multiplication modulo125
shows that $(a_1,w_1)$ are independent additive homomorphisms
and their joint kernel is precisely the subgroup for $T$.
Thus $T/S$ is Galois elementary abelian; its high line is $w_1=0$.
This verifies normality as well as counting an orbit.

For the quotient orientation the available comparison is the actual
image module $M_i[125]/H$ in $Q_i$, of type
$\mathbf Z/125\oplus\mathbf Z/25$. Take
\[
                   P^+=q(e_2/125),\qquad P=5P^+,
                                                        \tag{B-RD96.4}
\]
where $q:M_i\longrightarrow Q_i$ is the actual high isogeny.
Since $25(e_2/125)\notin H$, this point has exact order125.
Its $25$-multiple is the old first-level image and is the canonical
high point of $Q_i$. The original comparison identifies (B-RD96.4),
including its stabilizer, through the actual quotient image module.

The quotient lattice is $(e_1/5,e_2)$. Reorder it with its high
vector first, as $(e_2,e_1/5)$. Its matrix is
\[
                   \begin{pmatrix}1+5d&w\\5v&1+5a\end{pmatrix}.
                                                        \tag{B-RD96.5}
\]
Its entire reduction modulo125 is not needed. Its first column
modulo125 is determined by the original $125$-matrix. The stabilizer
of (B-RD96.4) is exactly $v_0=v_1=d_0=d_1=0$, and that of $P$
is $v_0=d_0=0$. The quotient residual field $R_o$ is the field
of $w_0$. Thus $S$ has stabilizer $w_0=v_0=d_0=0$, and $T$
adds the equations $d_1=v_1=0$. Their degrees are again those
in (B-RD96.1). On the subgroup for $S$, $(d_1,v_1)$ are
independent additive coordinates with the indicated kernel.
The high line is $v_1=0$. This construction has used only the
selected image point, not hypothetical extra quotient torsion.

In either orientation the two successive high inverse Newton
polygons are single segments, by the full coefficient estimate
of B49 and the actual quotient calculation of B69. More explicitly,
for parameter $5b>20$ the high first-level value is $5(24-b)$,
and successive inverse values are $(24-b)/5$ and $(24-b)/125$.
The degree-five coefficient vertex has value $20b$ and lies strictly
above the relevant degree-25 chords; all other vertices obey the
same complete coefficient estimate. This applies to the selected
point in the quotient orientation since its first-level target is
high. Multiplication by $[T:L]=3125$ gives $A_i$.
The low and high first-level values over $L$ are $b_i$ and
$5(24-b_i)$, giving $D_{0,i}$ and $D_{1,i}$. Formal translation
changes a coordinate by the translating coordinate times a unit.
Finally $v_L(5)=480$, giving the last equality of (B-RD96.2).
\end{proof}

\begin{lemma}[Opposite point fields have their high line last]
\label{lem:B-RD96-opposite}
\status{Candidate; complete proof awaiting independent audits.}
Consider an actual primitive low $25$-point field $Z$ in either
orientation, over that orientation's residual field $R_*$. It has
degree125 over $L$, is Galois of degree25 over $R_*$, and has the
two coordinate rows
\[
 A_i=5m_i,\qquad D_{0,i}=125m_i,\qquad
 D_{1,i}=15000-625m_i,\qquad v_Z(5)=60000.
                                                        \tag{B-RD96.6}
\]
For each of
\[
             (m_1,m_2)=(15,19),(11,19),(7,19),(9,17),(9,13),
                                                        \tag{B-RD96.7}
\]
its last ramification line is the common high-translation line.
Let $t'$ be the conductor of the diagonal character, which kills
that line, and $C'$ the off-diagonal conductor. Then
\[
\begin{array}{c|c|r}
 (m_1,m_2)&t'&C'\\\hline
 (15,19)&24\le t'\le336,\quad24\mid t'&606\\
 (11,19)&264&606\\
 (7,19)&168&606\\
 (9,17)&216&858\\
 (9,13)&216&1362.
\end{array}                                             \tag{B-RD96.8}
\]
\end{lemma}
\begin{proof}
The full group (B-RD96.3), reduced modulo25, gives a low-point
orbit of degree125. Over its residual field the translations
are the two independent additive second-column coordinates,
one diagonal and one high off-diagonal. The opposite quotient
matrix gives the same assertion. The low inverse value $m_i/25$
over $L$ proves (B-RD96.6).

The two translation axes have distinct tame characters, namely
the trivial character and a primitive fourth root in $\mathbf F_5$.
They cannot have just one ramification break. Otherwise the
injective leading-coefficient map of that ramification group
would identify tame conjugation with one scalar in the residue
field, incompatible with two different $\mathbf F_5$ eigenvalues.
There are therefore two lower breaks $t<u$. Denote their last
line by $T_*$. If both parameters are prime to five and $T_*$
is not the common high line, Lemma~\ref{lem:B-RD87-coefficient}
gives the same largest cyclic quotient conductor as both
$24m_1$ and $24m_2$. This is impossible for different parameters.

It remains to treat $(15,19)$ without any tracefree assumption.
Suppose again that the last line is not high. The prime row gives
\[
                           u+4t=120m_2.               \tag{B-RD96.9}
\]
For the other row $A=5m_1=75$ has five-adic order two. Use the
fixed coefficient expansion of B60 in a uniformizer of $Z$.
Let $N$ be the first nonzero exponent of five-adic order one,
and $J$ the first prime-to-five exponent, allowing either to be
infinite. Below $v_Z(5)$ the first possible change at break $h$ is
\[
                 \min\{A+25h,\ N+5h,\ J+h\}.
                                                        \tag{B-RD96.10}
\]
Higher exponents divisible by25 change later than the first term;
higher exponents of order one or zero change later than their
respective terms. If the minimum is unique it is the actual
change. All integers used here are below the stated cutoff.

An element of the high line has break $t$ and displacement $D_1$,
whereas a nonzero element of $T_*$ has break $u$ and displacement
$D_0$. Thus the minimum at $u$ is at most $D_0$, and its value
at $t$ is strictly less than $D_0<D_1$. At least two terms in
(B-RD96.10) must tie at $t$.
If the order-one and ordinary terms tie at the minimum, then
$J=N+4t$, and the ordinary term is uniquely first at $u$.
It would give $D_0=N+(u+4t)>A+120m_2>D_0$.
If the order-two and ordinary terms tie, the same reasoning gives
$J=A+24t$ and $D_0=A+20t+(u+4t)>D_0$.
These cases also cover a three-way tie. In the remaining case
$N=A+20t$ and $J>A+24t$. At $u$ the order-one term has value
$A+5(u+4t)=A+600m_2>D_0$, and the order-two term is even later.
Only the ordinary term could produce $D_0$, but
$J+u>A+20t+120m_2>D_0$. This exhausts the possibilities,
including absent coefficients, and proves that the high line is last.

Now the numerical proof of Lemma~\ref{lem:B-RD94-common-high}
applies to the actual rows (B-RD96.6): that proof uses only those
rows and the just-established last-high-line property. It uses
no tracefree relation. For two prime parameters it gives
$t'=24m_1$; using the larger row, which is strictly above this
critical break, then gives $C'=3000-126m_2$.
For $(15,19)$ the full mixed coefficient argument there gives
$0<t'<360$ and the same value $C'=606$. The diagonal cyclic
character is tame invariant. Its field has degree25 over $L$,
so a tame eigen-uniformizer has multiplier $\zeta$ because
$25\equiv1\pmod {24}$. The ramification leading map gives
$24\mid t'$, hence $24\le t'\le336$. This proves the table.
\end{proof}

\begin{proposition}[An actual large low-terminal conductor]
\label{prop:B-RD96-conductor}
\status{Candidate; complete proof awaiting independent audits.}
Orient the parameters as one of
\[
                    (5,9),(5,13),(5,17),(7,15),(11,15).
                                                        \tag{B-RD96.11}
\]
For the actual extension $T/S$ in Lemma~\ref{lem:B-RD96-selected},
let $z$ be the conductor of its cyclic quotient killing the high
translation line. Then the following strict bounds hold:
\[
\begin{array}{c|r|r|r}
 (b_1,b_2)&z>&9\min_i(D_{0,i}-A_i)/25&
                              \max_i(D_{0,i}-A_i)/5\\\hline
 (5,9)&7086&5454&5550\\
 (5,13)&8814&5454&8070\\
 (5,17)&11118&5454&10590\\
 (7,15)&16266&7722&9330\\
 (11,15)&28866&12258&9330.
\end{array}                                             \tag{B-RD96.12}
\]
\end{proposition}
\begin{proof}
First use the original orientation. Its opposite residual field
$R_*$ is the field of $w_0$. On the subgroup over $R_*$, the
three functions $v_0,a_0,w_1$ are independent additive characters.
The actual opposite low $25$-point field $Z/R_*$ has characters
$(a_0,w_1)$; by the preceding lemma their conductors are $t',C'$.
The actual extension $S/R_*$ has characters $(v_0,a_0)$.

Write $\sigma$ for the conductor of $v_0$ over $R_*$. Conjugating
by an opposite residual unipotent of parameter $t$ sends $a_0$
to $a_0\pm t v_0$. Conductor invariance therefore gives
$\sigma\le t'$. The two characters have different tame weights;
the same injective-leading-map argument used above rules out
equality. Consequently
\[
                              0<\sigma<t'<C'.          \tag{B-RD96.13}
\]
This uses the full $M_2$ congruence image, not a tracefree diagonal
identity. In particular the independence of all three characters
is preserved.

The actual compositum $SZ/R_*$ has these three independent
characters with strictly increasing conductors. Its upper
filtration therefore drops at $\sigma,t',C'$: a linear combination
of characters of different conductors has the largest conductor
occurring in it. The subgroup for $S$ is the last cyclic direction.
The exact Herbrand conversion, equivalently the quotient-average
formula in this elementary-abelian compositum, gives
\[
\begin{split}
 z&=\psi_{S/R_*}(C')
    =\sigma+5(t'-\sigma)+25(C'-t')\\
  &=25C'-20t'-4\sigma>25C'-24t'.
\end{split}                                             \tag{B-RD96.14}
\]
This is the conductor of the same character $w_1$ on $T/S$:
the actual common cyclic field is $S(w_1)=SZ$, and its kernel
in $T/S$ is the high line $w_1=0$. We have identified the actual
character and its base field before applying the conversion.

For the quotient orientation, use the actual selected point of
(B-RD96.4). The opposite residual field is now the field of $v_0$.
The three independent additive characters on its subgroup are
$w_0,d_0,v_1$. The opposite low-point field has characters
$(d_0,v_1)$, whereas $S$ has characters $(w_0,d_0)$.
The same conjugation gives the strict ordering
$\operatorname{cond}(w_0)<\operatorname{cond}(d_0)
 <\operatorname{cond}(v_1)$, and (B-RD96.14) computes the
actual $v_1$ quotient of $T/S$. Thus it applies equally to this
orientation without providing all quotient $125$-torsion.

The sorted opposite pairs of (B-RD96.11) are exactly (B-RD96.7)
in their listed order. Insert their $t',C'$ from (B-RD96.8)
into (B-RD96.14); in the first row use $t'\le336$.
The resulting strict lower bounds are the second column of
(B-RD96.12). The other columns follow directly from
$D_{0,i}-A_i=3150b_i-600$. This proves every table entry.
The conversion uses the actual-field conventions and the complete
proofs of Proposition3.7, Theorem3.10 and Proposition3.13(ii)
of \cite[pp.~14--16]{ColmezCorpsLocaux}, already retained and read
for the preceding ramification results.
\end{proof}

\begin{lemma}[A large quotient break forces equal displacement gaps]
\label{lem:B-RD96-gap}
\status{Candidate; complete proof awaiting independent audits.}
Let $T/S$ be an actual totally ramified elementary-abelian degree25
extension of
complete mixed-characteristic $(0,5)$ fields with perfect residue
field. Let $H$ be a line in its group. Suppose coordinates $x_i$
have valuations $A_i>0$, with $5\mid A_i$, and fixed displacements
$D_{1,i}$ on $H\setminus\{0\}$ and $D_{0,i}$ outside $H$, where
$A_i<D_{0,i}<D_{1,i}<v_T(5)$. Put $\Delta_i=D_{0,i}-A_i$.
If the conductor $z$ of the cyclic quotient killing $H$ satisfies
\[
 z>\frac9{25}\min_i\Delta_i,
 \qquad z>\frac15\max_i\Delta_i,
                                                        \tag{B-RD96.15}
\]
then all $D_{1,i}-D_{0,i}$ are equal.
\end{lemma}
\begin{proof}
Choose fixed residue representatives in $\mathcal O_S$ as in
B60, and expand each $x_i$ in a uniformizer of $T$. Up to the
stated cutoff, every exponent divisible by five has its first
possible change at degree at least $A_i+5h$ at a break $h$.
The first nonzero ordinary exponent $J_i$ changes at degree
$J_i+h$, with nonzero coefficient; all later ordinary terms
change later. This remains true if $A_i$ is divisible by25 or125.
Terms with five-divisible binomial coefficients lie beyond the
cutoff. No nonzero coefficient at any specified ordinary degree
is presumed.

If there is a single break, it equals $z$. The second inequality
puts all divisible-exponent changes after $D_{0,i}$. An element
outside $H$ therefore forces a finite first ordinary term with
$J_i+z=D_{0,i}$. Its leading coefficient is nonzero for every
nonidentity element at that break, including those in $H$.
This would give $D_{1,i}=D_{0,i}$, impossible.

Suppose there are two lower breaks $t<u$, with last line $T_1$.
If $T_1\ne H$, choose an index with minimal $\Delta_i$.
An element of $T_1$ is outside $H$, so the unconditional
fixed-coordinate inequality of B60 gives $u\le\Delta_i$.
If $t>\Delta_i/5$, all divisible changes are after $D_{0,i}$
already at $t$. The displacement at $u$ forces
$J_i+u=D_{0,i}$; at $t$ this ordinary term is uniquely first
and gives $J_i+t<D_{0,i}<D_{1,i}$ for an element of $H$.
This is impossible. Hence $t\le\Delta_i/5$, and the exact
cyclic-quotient average gives
\[
 z=\frac{u+4t}{5}\le\frac{9\Delta_i}{25},
                                                        \tag{B-RD96.16}
\]
contrary to (B-RD96.15).

The last line must therefore be $H$, and $t=z$. At the first
break the second inequality in (B-RD96.15) puts every divisible
term after $D_{0,i}$. Thus the outside-$H$ displacement forces
$J_i+t=D_{0,i}$. At the later break $u$, divisible terms have
increased at least five times as fast, whereas this ordinary
term has increased by $u-t$. It remains uniquely first below
the prescribed displacement cutoff, so
\[
 D_{1,i}=J_i+u,
 \qquad D_{1,i}-D_{0,i}=u-t.
                                                        \tag{B-RD96.17}
\]
If this putative first change were beyond the prescribed cutoff,
it could not produce $D_{1,i}$; thus the same equality is forced
in all cases. This proves the assertion, including every possible
vanishing of earlier ordinary coefficients.
\end{proof}

\begin{theorem}[The residual-wild comparison is impossible at the first wild index]
\label{thm:B-RD96-residual-wild}
\status{Candidate; complete proof awaiting independent audits.}
For $p=5,e_K=20$, no actual $125$-comparison of two whole
potentially supersingular height-two groups with different finite
special characters has nontrivial residual wild inertia.
\end{theorem}
\begin{proof}
B91--B93 exclude the non-both-large cases, and B94--B95 exclude
layer shape$(3,4)$. The remaining image is exactly (B-RD96.3).
The common high-line character and finite determinant give the
same finite parameter census as at the end of B94, independently
of the layer shape:
\[
\begin{gathered}
 (5,9),(5,13),(5,17),(7,15),(7,19),\\
 (9,13),(9,17),(11,15),(11,19),(15,19).
\end{gathered}                                          \tag{B-RD96.18}
\]
Indeed $b_i$ are odd, $b_1\equiv b_2\pmod4$, and the conditions
$c_i\equiv1\pmod4$, $c_i\equiv5b_i\pmod6$ and
$c_1-5b_1=c_2-5b_2\pmod {24}$ give this list after removing
equal unordered special orbits $\{c_i,5c_i\}$. The other two
same-congruence pairs $(7,11),(13,17)$ have equal special orbits.

Taking the sorted opposite pair $(24-b_2,24-b_1)$ groups these
ten rows into five pairs. Choose the orientation in (B-RD96.11).
Lemma~\ref{lem:B-RD96-selected} supplies the actual selected
point and common field even when this orientation is the high
quotient. Proposition~\ref{prop:B-RD96-conductor} verifies
both strict hypotheses of Lemma~\ref{lem:B-RD96-gap}.
Consequently $D_{1,1}-D_{0,1}=D_{1,2}-D_{0,2}$. But (B-RD96.2)
gives
\[
              D_{1,i}-D_{0,i}=375000-18750b_i,
                                                        \tag{B-RD96.19}
\]
which is different for $b_1<b_2$. This contradiction removes
every row of the full Iwahori case, and completes the residual-wild
exclusion. At no point has an extra torsion level been supplied
to the high quotients.
\end{proof}

\begin{remark}[Exact remaining scope and finite supplement]
\label{rem:B-RD96-scope}
\status{Candidate; scope accompanies the unpromoted theorem.}
If proved, this removes the remaining residual-wild branch rather
than realizing or excluding the irreducible-residual branch.
The latter still has the wild image orders $5^5$ through $5^8$
left by B81. The unrestricted least index at $p=5,e=20$ remains
in $\{2,3\}$. No existence of any remaining comparison is inferred.

The full B60 fixed-coefficient argument, B49/B69 actual inverse
Newton proof, B87 coefficient lemma, B91 residual reduction and
the complete B94/B95 proofs were reread for this argument.
No new external-source opening is claimed. The selected-point
construction, actual three-character compositum and general
coefficient lemma carry the proof; a finite arithmetic supplement
will be retained separately and will not replace these arguments.
\end{remark}
